\documentclass[a4paper,10pt]{article}
\usepackage[left=1in, right=1in, top=1in, bottom=1in]{geometry}
\usepackage{enumitem}
\usepackage{titlesec}
\usepackage{parskip}
\usepackage{xcolor}
\usepackage{fontawesome5}
\usepackage[hidelinks]{hyperref}
\usepackage{fancyhdr}
\usepackage{needspace}

% Define an olive green color
\definecolor{olivegreen}{RGB}{85,107,47}

\titleformat{\section}{\large\bfseries\color{olivegreen}}{}{0em}{}[\titlerule]
\titleformat{\subsection}[runin]{\bfseries\color{olivegreen}}{}{0em}{}[:]

\pagestyle{fancy}
\fancyhf{}
\fancyfoot[L]{Last updated: \today}

\begin{document}

\begin{center}
    \textbf{\Huge Antonio Neto} \\
    \vspace{2mm}
    \faIcon{briefcase} AI Engineer $|$ Computer Science Student
\end{center}

\vspace{4mm}
\begin{center}
    \faIcon{graduation-cap} \href{https://lattes.cnpq.br/6284651096866272}{Lattes} \hspace{1.5cm}
    \faIcon{github} \href{https://github.com/nietus}{Github} \hspace{1.5cm}
    \faIcon{envelope} \href{mailto:antonio@deeptera.com}{Email} \hspace{1.5cm}
    \faIcon{linkedin} \href{https://www.linkedin.com/in/antonioniet/}{Linkedin} \hspace{1.5cm}
    \faIcon{globe} \href{https://antonio.niet.com.br}{Website}
\end{center}

\space
\section*{\faIcon{user} Professional Profile}
I am a Brazilian AI Engineer and Computer Science student at PUC Minas, working with applied AI, optimization, infrastructure, and logistics problems. My research and development work covers LLM-based multi-agent systems, metaheuristics, and machine learning applied to industrial and logistics scenarios, with five accepted papers. I also study Mandarin at UFMG's Confucius Institute. I am highly focused, detail-oriented, and committed to everything I do.

\section*{\faIcon{building} Professional Experience}

\subsection*{Deeptera \hfill Apr 2026 -- Present}
\textit{Artificial Intelligence Engineer} \\
Part-time remote role focused on research and development
\begin{itemize}
    \item Research and development on optimization problems and multi-agent systems.
    \item Designed architectures for intelligent systems integrated into conventional software.
    \item Authored papers accepted at SBPO and SIMPEP on optimization and multi-agent systems.
\end{itemize}

\subsection*{Deeptera \hfill Aug 2024 -- Apr 2026}
\textit{Internship in Computing and Data} \\
Research and development
\begin{itemize}
    \item Published article on multi-agent systems approved at SBAI.
    \item Participated in projects funded by FAPEMIG.
    \item Designed database architecture and Django backend using Azure and Azure DevOps infrastructure.
    \item Provided optimization and backend services for VLI.
    \item Conducted specialized development and research related to multi-agent systems, with application of RAG algorithms, personalized tools and scalable services.
    \item Applied genetic algorithms to solve optimization problems.
\end{itemize}

\subsection*{Hackatruck \hfill 2023}
\textit{Project Participant} \\
Developed iOS applications as part of a month-long group project.
\begin{itemize}
    \item Collaborated with team members to design and implement app functionalities.
    \item Used Swift, Node Red and Azure for development, handling various APIs.
\end{itemize}

\section*{\faIcon{book} Publications}
\begin{itemize}[label=\faIcon{file-alt}]
    \needspace{4\baselineskip}\item \textbf{Uma Abordagem Baseada em Atenção para o Problema de Coletas e Entregas Dinâmicas.} A. Neto, R. D. Macedo, A. H. Duque, S. J. de Almeida. LVIII SBPO, 2026. \\
    Attention Model trained with Reinforcement Learning and combined with VNS for the Dynamic Pickup and Delivery Problem, evaluated on the ICAPS 2021 benchmark.
    \needspace{4\baselineskip}\item \textbf{DeepFactory: Sistema Multiagentes para Gestão e Otimização de Produção Industrial.} R. D. Macedo, A. Neto, A. H. Duque, S. J. de Almeida. LVIII SBPO, 2026. \\
    Stage-oriented multi-agent platform integrating MIP scheduling, demand forecasting, predictive maintenance and computer-vision quality control.
    \needspace{4\baselineskip}\item \textbf{Otimização do Sequenciamento de Operações de Carregamento e Descarregamento de Navios Graneleiros com Restrição de Trim via Algoritmo Genético.} A. Neto, R. D. Macedo, C. S. Almeida et al. (with VLI Multimodal). LVIII SBPO, 2026. \\
    Genetic algorithm with trim constraints validated on 112 real operational plans from a multimodal logistics company.
    \needspace{4\baselineskip}\item \textbf{Roteamento Multi-Período de Drones com Janelas de Tempo e Disponibilização Gradual de Estações de Recarga.} A. S. C. Neto, Z. K. G. do Patrocínio. SIMPEP, 2026. \\
    Undergraduate thesis: MP-EAVNS, an energy-aware Variable Neighborhood Search for a multi-period extension of the E-VRPTW.
    \needspace{4\baselineskip}\item \textbf{Proposta de um Sistema Multiagentes baseado em LLM: Integração em um Software de Gerenciamento e Otimização de Carregamento e Descarregamento de Navios.} A. Neto, C. S. Almeida, R. D. Macedo, A. H. Duque, S. J. de Almeida. SBAI/SBSE, 2025. \\
    Hierarchical LLM multi-agent system (Manager, Helper, Optimizer and Data Analytics) integrated into a conventional ship loading and unloading system.
\end{itemize}

\section*{\faIcon{graduation-cap} Education}

\subsection*{PUC Minas \hfill 2023 -- Present}
\textit{Bachelor of Computer Science} \\
I have worked with various technologies at PUC and specialized in machine learning, data science, and AI technologies outside of college. My undergraduate thesis addresses multi-period drone routing with time windows and gradual charging-station deployment, advised by Prof. Zenilton K. G. do Patrocínio.

\subsection*{Confucius Institute (UFMG) \hfill 2024 -- Present}
\textit{Mandarin Student (HSK 4)} \\
I study Mandarin out of a deep appreciation for the culture and its relevance to the field of AI.

\section*{\faIcon{cogs} Key Technical Skills}
\begin{center}
    \begin{itemize}[label=\faIcon{check}, itemsep=-3pt]
        \item \textbf{Programming Languages:} Python, C, TypeScript
        \item \textbf{Agent Systems:} LangChain, LangGraph, RAG, MCP
        \item \textbf{Machine Learning:} PyTorch, Scikit-learn, XGBoost
        \item \textbf{Data Science:} Pandas, NumPy, Matplotlib, Seaborn
        \item \textbf{Backend:} Django, FastAPI, REST APIs
        \item \textbf{Databases:} PostgreSQL, MySQL
        \item \textbf{DevOps:} Azure DevOps, Docker
        \item \textbf{Version Control:} Git, GitHub
        \item \textbf{Optimization:} AMPL, pyomo, GLPK/SCIP solvers, metaheuristics (GA, VNS, ALNS, GRASP)
        \item \textbf{Cloud:} Azure, Google Cloud

    \end{itemize}
\end{center}

\section*{\faIcon{certificate} Certifications}
\begin{itemize}[label=\faIcon{certificate}]
    \item Cloud Services Practices using Swift with emphasis on Cognitive Services
    \item Algorithms and OOP with Swift, JavaScript, and RESTful APIs
    \item Agriculture 4.0: Optimizing with Artificial Intelligence
\end{itemize}

\section*{\faIcon{language} Languages}
\begin{itemize}[label=\faIcon{globe}]
    \item Portuguese - Native
    \item English - Fluent
    \item Mandarin - HSK 4
\end{itemize}

\end{document}
